\documentclass[11pt]{article}
\usepackage[margin=1in]{geometry}
\usepackage{booktabs}
\usepackage{longtable}
\usepackage{hyperref}

\setlength{\emergencystretch}{3em}
\newcommand{\aal}[1]{\texttt{#1}}

\title{AAL90 to Yeo7 mapping}
\author{Juan Manuel Mnti}
\date{July 27, 2026}

\begin{document}
\maketitle

\begin{abstract}
Description of the mapping of AAL90 brain regions into the canonical Yeo 7-network
resting-state parcellation, with an optional subcortical extension for full 90-parcel
coverage.
\end{abstract}

\section*{Definition of the Yeo7 networks}

We parcellated the brain into 90 cortical and subcortical regions using the Automated
Anatomical Labeling (AAL) atlas in MNI space~\cite{tzourio2002,rolls2020}, and assigned
each cortical parcel to one of the seven canonical resting-state networks defined by
Yeo et al.~\cite{yeo2011} from clustering of resting-state functional connectivity in
1000 subjects. Four subcortical structures (Caudate, Putamen, Pallidum, Thalamus) have
no cortical Yeo7 assignment; we group them separately as a \emph{Subcortical extension}
so that analyses requiring full 90-parcel coverage are not forced to drop them. In
particular, we defined the following bilateral networks in terms of AAL90 regions (left
and right hemispheres):

\subsection*{Visual.}
The visual network comprised occipital cortex including primary and higher-order visual
areas~\cite{yeo2011}, namely bilateral \aal{Calcarine}, \aal{Cuneus}, \aal{Lingual},
\aal{Occipital\_Sup}, \aal{Occipital\_Mid}, \aal{Occipital\_Inf}, and \aal{Fusiform}.

\subsection*{Somatomotor.}
The somatomotor network comprised precentral, postcentral, and supplementary motor
regions, together with the opercular and auditory-adjacent parcels that load onto the
Yeo7 somatomotor component~\cite{yeo2011}. We included bilateral \aal{Precentral},
\aal{Rolandic\_Oper}, \aal{Supp\_Motor\_Area}, \aal{Postcentral},
\aal{Paracentral\_Lobule}, and \aal{Heschl}.

\subsection*{Dorsal attention.}
The dorsal attention network was defined as superior and inferior parietal regions,
approximating the intraparietal sulcus system described in attentional control
studies~\cite{corbetta2002}. We included bilateral \aal{Parietal\_Sup} and
\aal{Parietal\_Inf}.

\subsection*{Ventral attention / Salience.}
The ventral attention / salience network comprised anterior insula, cingulate, and
opercular regions, consistent with insula--anterior cingulate models of the salience
system~\cite{seeley2007}, together with the superior temporal parcels that fall in the
Yeo7 ventral attention component. We included bilateral \aal{Insula},
\aal{Cingulum\_Ant}, \aal{Cingulum\_Mid}, \aal{SupraMarginal}, and \aal{Temporal\_Sup}.

\subsection*{Limbic.}
The limbic network comprised orbitofrontal, temporal pole, and medial temporal regions,
in line with the Yeo7 limbic component~\cite{yeo2011}. We included bilateral
\aal{Frontal\_Sup\_Orb}, \aal{Frontal\_Mid\_Orb}, \aal{Frontal\_Inf\_Orb},
\aal{Olfactory}, \aal{Frontal\_Med\_Orb}, \aal{Rectus}, \aal{Amygdala},
\aal{Temporal\_Pole\_Sup}, and \aal{Temporal\_Pole\_Mid}.

\subsection*{Frontal control.}
The frontal control (frontoparietal) network was defined as dorsolateral and inferior
frontal regions~\cite{yeo2011}, including bilateral \aal{Frontal\_Sup},
\aal{Frontal\_Mid}, \aal{Frontal\_Inf\_Oper}, and \aal{Frontal\_Inf\_Tri}. These parcels
approximate dorsolateral prefrontal cortex within the AAL scheme.

\subsection*{Default mode.}
The default mode network comprised medial prefrontal, posterior cingulate/precuneus,
inferior parietal, and medial temporal regions, in line with canonical DMN
descriptions~\cite{raichle2001,yeo2011}. We included bilateral
\aal{Frontal\_Sup\_Medial}, \aal{Cingulum\_Post}, \aal{Hippocampus},
\aal{ParaHippocampal}, \aal{Angular}, \aal{Precuneus}, \aal{Temporal\_Mid}, and
\aal{Temporal\_Inf}.

\subsection*{Subcortical extension (non-Yeo7).}
Caudate, Putamen, Pallidum, and Thalamus are not part of the cortical Yeo7 parcellation.
When full 90-parcel coverage is required we group bilateral \aal{Caudate},
\aal{Putamen}, \aal{Pallidum}, and \aal{Thalamus} into a separate \emph{Subcortical
extension} network rather than force them into a cortical network they do not belong
to.

\section*{Reference tables}

For transparency and reproducibility, the full list of AAL90 parcels belonging to each
network (and their indices in the AAL-90 ordering) is provided in
Table~\ref{tab:yeo7}.

\begin{longtable}{lll}
\caption{Assignment of AAL-90 parcels to Yeo7 networks. Indices refer to the 1-based
position in the AAL-90 non-symmetric ordering (Left then Right).}
\label{tab:yeo7}\\
\toprule
Network & AAL label & AAL index \\
\midrule
\endfirsthead
\toprule
Network & AAL label & AAL index \\
\midrule
\endhead
\midrule
\multicolumn{3}{r}{\emph{continued on next page}} \\
\endfoot
\bottomrule
\endlastfoot
Default Mode & Frontal\_Sup\_Medial\_L & 12 \\
 & Frontal\_Sup\_Medial\_R & 57 \\
 & Cingulum\_Post\_L & 18 \\
 & Cingulum\_Post\_R & 63 \\
 & Hippocampus\_L & 19 \\
 & Hippocampus\_R & 64 \\
 & ParaHippocampal\_L & 20 \\
 & ParaHippocampal\_R & 65 \\
 & Angular\_L & 33 \\
 & Angular\_R & 78 \\
 & Precuneus\_L & 34 \\
 & Precuneus\_R & 79 \\
 & Temporal\_Mid\_L & 43 \\
 & Temporal\_Mid\_R & 88 \\
 & Temporal\_Inf\_L & 45 \\
 & Temporal\_Inf\_R & 90 \\
\midrule
Dorsal Attention & Parietal\_Sup\_L & 30 \\
 & Parietal\_Sup\_R & 75 \\
 & Parietal\_Inf\_L & 31 \\
 & Parietal\_Inf\_R & 76 \\
\midrule
Ventral Attention / Salience & Insula\_L & 15 \\
 & Insula\_R & 60 \\
 & Cingulum\_Ant\_L & 16 \\
 & Cingulum\_Ant\_R & 61 \\
 & Cingulum\_Mid\_L & 17 \\
 & Cingulum\_Mid\_R & 62 \\
 & SupraMarginal\_L & 32 \\
 & SupraMarginal\_R & 77 \\
 & Temporal\_Sup\_L & 41 \\
 & Temporal\_Sup\_R & 86 \\
\midrule
Frontal Control & Frontal\_Sup\_L & 2 \\
 & Frontal\_Sup\_R & 47 \\
 & Frontal\_Mid\_L & 4 \\
 & Frontal\_Mid\_R & 49 \\
 & Frontal\_Inf\_Oper\_L & 6 \\
 & Frontal\_Inf\_Oper\_R & 51 \\
 & Frontal\_Inf\_Tri\_L & 7 \\
 & Frontal\_Inf\_Tri\_R & 52 \\
\midrule
Limbic & Frontal\_Sup\_Orb\_L & 3 \\
 & Frontal\_Sup\_Orb\_R & 48 \\
 & Frontal\_Mid\_Orb\_L & 5 \\
 & Frontal\_Mid\_Orb\_R & 50 \\
 & Frontal\_Inf\_Orb\_L & 8 \\
 & Frontal\_Inf\_Orb\_R & 53 \\
 & Olfactory\_L & 11 \\
 & Olfactory\_R & 56 \\
 & Frontal\_Med\_Orb\_L & 13 \\
 & Frontal\_Med\_Orb\_R & 58 \\
 & Rectus\_L & 14 \\
 & Rectus\_R & 59 \\
 & Amygdala\_L & 21 \\
 & Amygdala\_R & 66 \\
 & Temporal\_Pole\_Sup\_L & 42 \\
 & Temporal\_Pole\_Sup\_R & 87 \\
 & Temporal\_Pole\_Mid\_L & 44 \\
 & Temporal\_Pole\_Mid\_R & 89 \\
\midrule
Somatomotor & Precentral\_L & 1 \\
 & Precentral\_R & 46 \\
 & Rolandic\_Oper\_L & 9 \\
 & Rolandic\_Oper\_R & 54 \\
 & Supp\_Motor\_Area\_L & 10 \\
 & Supp\_Motor\_Area\_R & 55 \\
 & Postcentral\_L & 29 \\
 & Postcentral\_R & 74 \\
 & Paracentral\_Lobule\_L & 35 \\
 & Paracentral\_Lobule\_R & 80 \\
 & Heschl\_L & 40 \\
 & Heschl\_R & 85 \\
\midrule
Visual & Calcarine\_L & 22 \\
 & Calcarine\_R & 67 \\
 & Cuneus\_L & 23 \\
 & Cuneus\_R & 68 \\
 & Lingual\_L & 24 \\
 & Lingual\_R & 69 \\
 & Occipital\_Sup\_L & 25 \\
 & Occipital\_Sup\_R & 70 \\
 & Occipital\_Mid\_L & 26 \\
 & Occipital\_Mid\_R & 71 \\
 & Occipital\_Inf\_L & 27 \\
 & Occipital\_Inf\_R & 72 \\
 & Fusiform\_L & 28 \\
 & Fusiform\_R & 73 \\
\midrule
Subcortical extension & Caudate\_L & 36 \\
 & Caudate\_R & 81 \\
 & Putamen\_L & 37 \\
 & Putamen\_R & 82 \\
 & Pallidum\_L & 38 \\
 & Pallidum\_R & 83 \\
 & Thalamus\_L & 39 \\
 & Thalamus\_R & 84 \\
\end{longtable}

\begin{thebibliography}{9}

\bibitem{yeo2011}
B. T. Thomas Yeo, Fenna M. Krienen, Jorge Sepulcre, Mert R. Sabuncu, Danial Lashkari,
Marisa Hollinshead, Joshua L. Roffman, Jordan W. Smoller, Lilla Z\"{o}llei, Jonathan R.
Polimeni, Bruce Fischl, Hesheng Liu, and Randy L. Buckner. The organization of the human
cerebral cortex estimated by intrinsic functional connectivity. \emph{Journal of
Neurophysiology}, 106(3):1125--1165, 2011.

\bibitem{tzourio2002}
N. Tzourio-Mazoyer, B. Landeau, D. Papathanassiou, F. Crivello, O. Etard, N. Delcroix,
B. Mazoyer, and M. Joliot. Automated anatomical labeling of activations in SPM using a
macroscopic anatomical parcellation of the MNI MRI single-subject brain.
\emph{NeuroImage}, 15(1):273--289, 2002.

\bibitem{rolls2020}
Edmund T. Rolls, Ching-Chang Huang, Chia-Chien Lin, Jian Feng, and Marc Joliot.
Automated anatomical labelling atlas 3. \emph{NeuroImage}, 206:116189, 2020.

\bibitem{raichle2001}
Marcus E. Raichle, Ann M. MacLeod, Abraham Z. Snyder, William J. Powers, Debra A.
Gusnard, and Gordon L. Shulman. A default mode of brain function. \emph{Proceedings of
the National Academy of Sciences of the United States of America}, 98(2):676--682,
2001.

\bibitem{corbetta2002}
Maurizio Corbetta and Gordon L. Shulman. Control of goal-directed and stimulus-driven
attention in the brain. \emph{Nature Reviews Neuroscience}, 3(3):201--215, 2002.

\bibitem{seeley2007}
William W. Seeley, Vinod Menon, Alan F. Schatzberg, Jeffrey Keller, Gary H. Glover,
Hilary Kenna, Allan L. Reiss, and Michael D. Greicius. Dissociable intrinsic
connectivity networks for salience processing and executive control. \emph{Journal of
Neuroscience}, 27(9):2349--2356, 2007.

\end{thebibliography}

\end{document}
